% Q03_problems.tex — the written problems of Q03 (First syntax II: modes,
% logic, tooling, trust). Built by `make -C paper paper` with Q03_answers.tex.
% B1 and B2 are the S8 lab (weakest preconditions and verification conditions
% by hand, 20 minutes); B3 (a blueprint) and B4 (the paper card) are the
% module's written pair (15 minutes).
\documentclass[11pt]{article}
% preamble.tex — shared preamble for the dafny-verus-learn-0 written track.
%
% Usage (every <ID>_problems.tex, <ID>_answers.tex and solutions/<ID>_solutions.tex):
%
%   \documentclass[11pt]{article}
%   % preamble.tex — shared preamble for the dafny-verus-learn-0 written track.
%
% Usage (every <ID>_problems.tex, <ID>_answers.tex and solutions/<ID>_solutions.tex):
%
%   \documentclass[11pt]{article}
%   % preamble.tex — shared preamble for the dafny-verus-learn-0 written track.
%
% Usage (every <ID>_problems.tex, <ID>_answers.tex and solutions/<ID>_solutions.tex):
%
%   \documentclass[11pt]{article}
%   \input{preamble}          % the Makefile always runs latexmk from paper/
%   \begin{document}
%   ...
%   \end{document}
%
% Engine: pdflatex. Dafny and Verus sources are mostly ASCII; the `literate`
% tables below map the common Unicode symbols to math glyphs inside listings.
% For other Unicode, extend a table, or build that file with lualatex
% (`latexmk -lualatex <file>`) after replacing the inputenc/fontenc lines with
% \usepackage{fontspec} and a monospace font with good coverage.
%
% Contents: program-logic macros (Hoare triples, wp, VCs, substitution,
% semantics brackets, typing judgments), inference rules (mathpartir),
% theorem and problem environments, the answer environment, the grader
% margin-note macro, and listings languages `dafny` and `verus` (Python uses
% listings' built-in `Python`), with inline forms \dfy{...}, \vrs{...}, \py{...}.
\usepackage[utf8]{inputenc}
\usepackage[T1]{fontenc}

\usepackage{amsmath}
\usepackage{amssymb}
\usepackage{amsthm}
\usepackage{stmaryrd}         % \llbracket \rrbracket for semantics brackets
\usepackage{mathpartir}       % inference rules: \inferrule, the mathpar environment
\usepackage{listings}
\usepackage{xcolor}
\usepackage[margin=1.1in]{geometry}

% ---------------------------------------------------------------------------
% Program logic.
%   \hoare{P}{c}{Q}       {P} c {Q}           (partial correctness)
%   \thoare{P}{c}{Q}      [P] c [Q]           (total correctness)
%   \provable{P}{c}{Q}    |- {P} c {Q}        (derivable in the proof system)
%   \WP{c}{Q}             wp(c, Q)            (\wp is LaTeX's Weierstrass p, so
%   \WLP{c}{Q}            wlp(c, Q)            the macro names are upper case)
%   \VC{c}{Q}             VC(c, Q)
%   \subst{Q}{x}{e}       Q[e/x]              (e replaces x in Q)
%   \sem{e}               [[e]]               (denotation)
%   \assign, \seqc        :=  and the sequencing semicolon
%   \impl                 ==> as a math arrow (LaTeX's own \iff gives <==>)
% ---------------------------------------------------------------------------
\newcommand{\hoare}[3]{\{\,#1\,\}\;#2\;\{\,#3\,\}}
\newcommand{\thoare}[3]{[\,#1\,]\;#2\;[\,#3\,]}
\newcommand{\provable}[3]{\vdash \hoare{#1}{#2}{#3}}
\newcommand{\WP}[2]{\mathrm{wp}(#1,\,#2)}
\newcommand{\WLP}[2]{\mathrm{wlp}(#1,\,#2)}
\newcommand{\VC}[2]{\mathrm{VC}(#1,\,#2)}
\newcommand{\subst}[3]{#1[#3/#2]}
\newcommand{\sem}[1]{\llbracket #1 \rrbracket}
\newcommand{\assign}{\mathrel{:=}}
\newcommand{\seqc}{\mathbin{;}\,}
\newcommand{\impl}{\Longrightarrow}

% ---------------------------------------------------------------------------
% Judgments.  \judgment{\Gamma}{e}{\tau}  ==>  Gamma |- e : tau
% ---------------------------------------------------------------------------
\newcommand{\judgment}[3]{#1 \vdash #2 : #3}
\newcommand{\judg}[3]{\judgment{#1}{#2}{#3}}          % short alias
\newcommand{\emptyjudgment}[2]{\cdot \vdash #1 : #2}  % empty context
\newcommand{\ctx}{\Gamma}
\newcommand{\turnstile}{\vdash}

% ---------------------------------------------------------------------------
% Theorem environments.
% ---------------------------------------------------------------------------
\theoremstyle{plain}
\newtheorem{theorem}{Theorem}[section]
\newtheorem{lemma}[theorem]{Lemma}
\newtheorem{corollary}[theorem]{Corollary}
\newtheorem{proposition}[theorem]{Proposition}
\theoremstyle{definition}
\newtheorem{definition}[theorem]{Definition}
% Written-track problems are numbered B1, B2, ... on both tracks (P<n> is a
% predict item). Label each one \label{prob:B<n>} and its answer \label{ans:B<n>}.
\newtheorem{problem}{Problem}
\renewcommand{\theproblem}{B\arabic{problem}}
\theoremstyle{remark}
\newtheorem*{remark}{Remark}
\newenvironment{answer}
  {\par\medskip\noindent\textbf{Answer.}\begin{quotation}\noindent}
  {\end{quotation}\medskip}

% ---------------------------------------------------------------------------
% Grader margin notes. Rendered form of the `% GRADER[...][...][...]:` comment
% convention (see template_answers.tex). Use inline where a visible note helps:
%   \gradernote{B2/induction-soundness/2: the invariant is not preserved by the last assignment}
% ---------------------------------------------------------------------------
\definecolor{gradercolor}{RGB}{160, 30, 30}
\newcommand{\gradernote}[1]{\textcolor{gradercolor}{\small\sffamily [GRADER: #1]}}

% ---------------------------------------------------------------------------
% Code listings: Dafny, Verus (Rust inside verus!{...}), Python.
% ---------------------------------------------------------------------------
\definecolor{codekeyword}{RGB}{0, 60, 130}
\definecolor{codecomment}{RGB}{90, 110, 90}
\definecolor{codestring}{RGB}{130, 60, 10}
\definecolor{codehole}{RGB}{180, 0, 0}

\lstdefinelanguage{dafny}{
  sensitive=true,
  morekeywords={
    abstract, allocated, as, assert, assume, bool, break, by, calc, case, char,
    class, codatatype, const, constructor, continue, datatype, decreases, else,
    ensures, exists, expect, export, extends, false, forall, fresh, function,
    ghost, if, import, in, include, int, invariant, is, iterator, label, lemma,
    map, match, method, modifies, modify, module, multiset, nat, new, newtype,
    null, old, opaque, opened, predicate, print, provides, reads, real, refines,
    requires, return, returns, reveal, reveals, seq, set, static, string, then,
    this, trait, true, twostate, type, unchanged, var, while, witness},
  emph={assume},
  emphstyle=\color{codehole}\bfseries,
  comment=[l]{//},
  morecomment=[s]{/*}{*/},
  morestring=[b]",
}

\lstdefinelanguage{verus}{
  sensitive=true,
  morekeywords={
    as, break, const, continue, crate, dyn, else, enum, false, fn, for, if,
    impl, in, let, loop, match, mod, move, mut, pub, ref, return, self, Self,
    static, struct, super, trait, true, type, unsafe, use, where, while,
    spec, proof, exec, open, closed, tracked, ghost, requires, ensures,
    recommends, invariant, decreases, assert, forall, exists, choose, by,
    broadcast, uninterp, axiom, returns, int, nat, Seq, Set, Map, Ghost, Tracked},
  emph={assume, admit},
  emphstyle=\color{codehole}\bfseries,
  comment=[l]{//},
  morecomment=[s]{/*}{*/},
  morestring=[b]",
}

\lstset{
  basicstyle=\ttfamily\small,
  keywordstyle=\color{codekeyword}\bfseries,
  commentstyle=\color{codecomment}\itshape,
  stringstyle=\color{codestring},
  showstringspaces=false,
  columns=fullflexible,
  keepspaces=true,
  breaklines=true,
  frame=leftline,
  framerule=0.6pt,
  xleftmargin=1em,
  aboveskip=0.8em,
  belowskip=0.8em,
  literate=
    {→}{{$\to$}}1        {←}{{$\leftarrow$}}1  {⇒}{{$\Rightarrow$}}1
    {⇔}{{$\Leftrightarrow$}}1                  {↦}{{$\mapsto$}}1
    {∀}{{$\forall$}}1    {∃}{{$\exists$}}1     {λ}{{$\lambda$}}1
    {∧}{{$\wedge$}}1     {∨}{{$\vee$}}1        {¬}{{$\neg$}}1
    {≠}{{$\neq$}}1       {≤}{{$\leq$}}1        {≥}{{$\geq$}}1
    {∈}{{$\in$}}1        {∉}{{$\notin$}}1      {⊆}{{$\subseteq$}}1
    {×}{{$\times$}}1     {·}{{$\cdot$}}1       {∅}{{$\varnothing$}}1
    {ℕ}{{$\mathbb{N}$}}1 {ℤ}{{$\mathbb{Z}$}}1
    {⊢}{{$\vdash$}}1     {⟦}{{$\llbracket$}}1  {⟧}{{$\rrbracket$}}1,
}

% Code environments and inline code.
\lstnewenvironment{dafnycode}{\lstset{language=dafny}}{}
\lstnewenvironment{veruscode}{\lstset{language=verus}}{}
\lstnewenvironment{pycode}{\lstset{language=Python}}{}
\newcommand{\dfy}[1]{\lstinline[language=dafny]{#1}}
\newcommand{\vrs}[1]{\lstinline[language=verus]{#1}}
\newcommand{\py}[1]{\lstinline[language=Python]{#1}}
          % the Makefile always runs latexmk from paper/
%   \begin{document}
%   ...
%   \end{document}
%
% Engine: pdflatex. Dafny and Verus sources are mostly ASCII; the `literate`
% tables below map the common Unicode symbols to math glyphs inside listings.
% For other Unicode, extend a table, or build that file with lualatex
% (`latexmk -lualatex <file>`) after replacing the inputenc/fontenc lines with
% \usepackage{fontspec} and a monospace font with good coverage.
%
% Contents: program-logic macros (Hoare triples, wp, VCs, substitution,
% semantics brackets, typing judgments), inference rules (mathpartir),
% theorem and problem environments, the answer environment, the grader
% margin-note macro, and listings languages `dafny` and `verus` (Python uses
% listings' built-in `Python`), with inline forms \dfy{...}, \vrs{...}, \py{...}.
\usepackage[utf8]{inputenc}
\usepackage[T1]{fontenc}

\usepackage{amsmath}
\usepackage{amssymb}
\usepackage{amsthm}
\usepackage{stmaryrd}         % \llbracket \rrbracket for semantics brackets
\usepackage{mathpartir}       % inference rules: \inferrule, the mathpar environment
\usepackage{listings}
\usepackage{xcolor}
\usepackage[margin=1.1in]{geometry}

% ---------------------------------------------------------------------------
% Program logic.
%   \hoare{P}{c}{Q}       {P} c {Q}           (partial correctness)
%   \thoare{P}{c}{Q}      [P] c [Q]           (total correctness)
%   \provable{P}{c}{Q}    |- {P} c {Q}        (derivable in the proof system)
%   \WP{c}{Q}             wp(c, Q)            (\wp is LaTeX's Weierstrass p, so
%   \WLP{c}{Q}            wlp(c, Q)            the macro names are upper case)
%   \VC{c}{Q}             VC(c, Q)
%   \subst{Q}{x}{e}       Q[e/x]              (e replaces x in Q)
%   \sem{e}               [[e]]               (denotation)
%   \assign, \seqc        :=  and the sequencing semicolon
%   \impl                 ==> as a math arrow (LaTeX's own \iff gives <==>)
% ---------------------------------------------------------------------------
\newcommand{\hoare}[3]{\{\,#1\,\}\;#2\;\{\,#3\,\}}
\newcommand{\thoare}[3]{[\,#1\,]\;#2\;[\,#3\,]}
\newcommand{\provable}[3]{\vdash \hoare{#1}{#2}{#3}}
\newcommand{\WP}[2]{\mathrm{wp}(#1,\,#2)}
\newcommand{\WLP}[2]{\mathrm{wlp}(#1,\,#2)}
\newcommand{\VC}[2]{\mathrm{VC}(#1,\,#2)}
\newcommand{\subst}[3]{#1[#3/#2]}
\newcommand{\sem}[1]{\llbracket #1 \rrbracket}
\newcommand{\assign}{\mathrel{:=}}
\newcommand{\seqc}{\mathbin{;}\,}
\newcommand{\impl}{\Longrightarrow}

% ---------------------------------------------------------------------------
% Judgments.  \judgment{\Gamma}{e}{\tau}  ==>  Gamma |- e : tau
% ---------------------------------------------------------------------------
\newcommand{\judgment}[3]{#1 \vdash #2 : #3}
\newcommand{\judg}[3]{\judgment{#1}{#2}{#3}}          % short alias
\newcommand{\emptyjudgment}[2]{\cdot \vdash #1 : #2}  % empty context
\newcommand{\ctx}{\Gamma}
\newcommand{\turnstile}{\vdash}

% ---------------------------------------------------------------------------
% Theorem environments.
% ---------------------------------------------------------------------------
\theoremstyle{plain}
\newtheorem{theorem}{Theorem}[section]
\newtheorem{lemma}[theorem]{Lemma}
\newtheorem{corollary}[theorem]{Corollary}
\newtheorem{proposition}[theorem]{Proposition}
\theoremstyle{definition}
\newtheorem{definition}[theorem]{Definition}
% Written-track problems are numbered B1, B2, ... on both tracks (P<n> is a
% predict item). Label each one \label{prob:B<n>} and its answer \label{ans:B<n>}.
\newtheorem{problem}{Problem}
\renewcommand{\theproblem}{B\arabic{problem}}
\theoremstyle{remark}
\newtheorem*{remark}{Remark}
\newenvironment{answer}
  {\par\medskip\noindent\textbf{Answer.}\begin{quotation}\noindent}
  {\end{quotation}\medskip}

% ---------------------------------------------------------------------------
% Grader margin notes. Rendered form of the `% GRADER[...][...][...]:` comment
% convention (see template_answers.tex). Use inline where a visible note helps:
%   \gradernote{B2/induction-soundness/2: the invariant is not preserved by the last assignment}
% ---------------------------------------------------------------------------
\definecolor{gradercolor}{RGB}{160, 30, 30}
\newcommand{\gradernote}[1]{\textcolor{gradercolor}{\small\sffamily [GRADER: #1]}}

% ---------------------------------------------------------------------------
% Code listings: Dafny, Verus (Rust inside verus!{...}), Python.
% ---------------------------------------------------------------------------
\definecolor{codekeyword}{RGB}{0, 60, 130}
\definecolor{codecomment}{RGB}{90, 110, 90}
\definecolor{codestring}{RGB}{130, 60, 10}
\definecolor{codehole}{RGB}{180, 0, 0}

\lstdefinelanguage{dafny}{
  sensitive=true,
  morekeywords={
    abstract, allocated, as, assert, assume, bool, break, by, calc, case, char,
    class, codatatype, const, constructor, continue, datatype, decreases, else,
    ensures, exists, expect, export, extends, false, forall, fresh, function,
    ghost, if, import, in, include, int, invariant, is, iterator, label, lemma,
    map, match, method, modifies, modify, module, multiset, nat, new, newtype,
    null, old, opaque, opened, predicate, print, provides, reads, real, refines,
    requires, return, returns, reveal, reveals, seq, set, static, string, then,
    this, trait, true, twostate, type, unchanged, var, while, witness},
  emph={assume},
  emphstyle=\color{codehole}\bfseries,
  comment=[l]{//},
  morecomment=[s]{/*}{*/},
  morestring=[b]",
}

\lstdefinelanguage{verus}{
  sensitive=true,
  morekeywords={
    as, break, const, continue, crate, dyn, else, enum, false, fn, for, if,
    impl, in, let, loop, match, mod, move, mut, pub, ref, return, self, Self,
    static, struct, super, trait, true, type, unsafe, use, where, while,
    spec, proof, exec, open, closed, tracked, ghost, requires, ensures,
    recommends, invariant, decreases, assert, forall, exists, choose, by,
    broadcast, uninterp, axiom, returns, int, nat, Seq, Set, Map, Ghost, Tracked},
  emph={assume, admit},
  emphstyle=\color{codehole}\bfseries,
  comment=[l]{//},
  morecomment=[s]{/*}{*/},
  morestring=[b]",
}

\lstset{
  basicstyle=\ttfamily\small,
  keywordstyle=\color{codekeyword}\bfseries,
  commentstyle=\color{codecomment}\itshape,
  stringstyle=\color{codestring},
  showstringspaces=false,
  columns=fullflexible,
  keepspaces=true,
  breaklines=true,
  frame=leftline,
  framerule=0.6pt,
  xleftmargin=1em,
  aboveskip=0.8em,
  belowskip=0.8em,
  literate=
    {→}{{$\to$}}1        {←}{{$\leftarrow$}}1  {⇒}{{$\Rightarrow$}}1
    {⇔}{{$\Leftrightarrow$}}1                  {↦}{{$\mapsto$}}1
    {∀}{{$\forall$}}1    {∃}{{$\exists$}}1     {λ}{{$\lambda$}}1
    {∧}{{$\wedge$}}1     {∨}{{$\vee$}}1        {¬}{{$\neg$}}1
    {≠}{{$\neq$}}1       {≤}{{$\leq$}}1        {≥}{{$\geq$}}1
    {∈}{{$\in$}}1        {∉}{{$\notin$}}1      {⊆}{{$\subseteq$}}1
    {×}{{$\times$}}1     {·}{{$\cdot$}}1       {∅}{{$\varnothing$}}1
    {ℕ}{{$\mathbb{N}$}}1 {ℤ}{{$\mathbb{Z}$}}1
    {⊢}{{$\vdash$}}1     {⟦}{{$\llbracket$}}1  {⟧}{{$\rrbracket$}}1,
}

% Code environments and inline code.
\lstnewenvironment{dafnycode}{\lstset{language=dafny}}{}
\lstnewenvironment{veruscode}{\lstset{language=verus}}{}
\lstnewenvironment{pycode}{\lstset{language=Python}}{}
\newcommand{\dfy}[1]{\lstinline[language=dafny]{#1}}
\newcommand{\vrs}[1]{\lstinline[language=verus]{#1}}
\newcommand{\py}[1]{\lstinline[language=Python]{#1}}
          % the Makefile always runs latexmk from paper/
%   \begin{document}
%   ...
%   \end{document}
%
% Engine: pdflatex. Dafny and Verus sources are mostly ASCII; the `literate`
% tables below map the common Unicode symbols to math glyphs inside listings.
% For other Unicode, extend a table, or build that file with lualatex
% (`latexmk -lualatex <file>`) after replacing the inputenc/fontenc lines with
% \usepackage{fontspec} and a monospace font with good coverage.
%
% Contents: program-logic macros (Hoare triples, wp, VCs, substitution,
% semantics brackets, typing judgments), inference rules (mathpartir),
% theorem and problem environments, the answer environment, the grader
% margin-note macro, and listings languages `dafny` and `verus` (Python uses
% listings' built-in `Python`), with inline forms \dfy{...}, \vrs{...}, \py{...}.
\usepackage[utf8]{inputenc}
\usepackage[T1]{fontenc}

\usepackage{amsmath}
\usepackage{amssymb}
\usepackage{amsthm}
\usepackage{stmaryrd}         % \llbracket \rrbracket for semantics brackets
\usepackage{mathpartir}       % inference rules: \inferrule, the mathpar environment
\usepackage{listings}
\usepackage{xcolor}
\usepackage[margin=1.1in]{geometry}

% ---------------------------------------------------------------------------
% Program logic.
%   \hoare{P}{c}{Q}       {P} c {Q}           (partial correctness)
%   \thoare{P}{c}{Q}      [P] c [Q]           (total correctness)
%   \provable{P}{c}{Q}    |- {P} c {Q}        (derivable in the proof system)
%   \WP{c}{Q}             wp(c, Q)            (\wp is LaTeX's Weierstrass p, so
%   \WLP{c}{Q}            wlp(c, Q)            the macro names are upper case)
%   \VC{c}{Q}             VC(c, Q)
%   \subst{Q}{x}{e}       Q[e/x]              (e replaces x in Q)
%   \sem{e}               [[e]]               (denotation)
%   \assign, \seqc        :=  and the sequencing semicolon
%   \impl                 ==> as a math arrow (LaTeX's own \iff gives <==>)
% ---------------------------------------------------------------------------
\newcommand{\hoare}[3]{\{\,#1\,\}\;#2\;\{\,#3\,\}}
\newcommand{\thoare}[3]{[\,#1\,]\;#2\;[\,#3\,]}
\newcommand{\provable}[3]{\vdash \hoare{#1}{#2}{#3}}
\newcommand{\WP}[2]{\mathrm{wp}(#1,\,#2)}
\newcommand{\WLP}[2]{\mathrm{wlp}(#1,\,#2)}
\newcommand{\VC}[2]{\mathrm{VC}(#1,\,#2)}
\newcommand{\subst}[3]{#1[#3/#2]}
\newcommand{\sem}[1]{\llbracket #1 \rrbracket}
\newcommand{\assign}{\mathrel{:=}}
\newcommand{\seqc}{\mathbin{;}\,}
\newcommand{\impl}{\Longrightarrow}

% ---------------------------------------------------------------------------
% Judgments.  \judgment{\Gamma}{e}{\tau}  ==>  Gamma |- e : tau
% ---------------------------------------------------------------------------
\newcommand{\judgment}[3]{#1 \vdash #2 : #3}
\newcommand{\judg}[3]{\judgment{#1}{#2}{#3}}          % short alias
\newcommand{\emptyjudgment}[2]{\cdot \vdash #1 : #2}  % empty context
\newcommand{\ctx}{\Gamma}
\newcommand{\turnstile}{\vdash}

% ---------------------------------------------------------------------------
% Theorem environments.
% ---------------------------------------------------------------------------
\theoremstyle{plain}
\newtheorem{theorem}{Theorem}[section]
\newtheorem{lemma}[theorem]{Lemma}
\newtheorem{corollary}[theorem]{Corollary}
\newtheorem{proposition}[theorem]{Proposition}
\theoremstyle{definition}
\newtheorem{definition}[theorem]{Definition}
% Written-track problems are numbered B1, B2, ... on both tracks (P<n> is a
% predict item). Label each one \label{prob:B<n>} and its answer \label{ans:B<n>}.
\newtheorem{problem}{Problem}
\renewcommand{\theproblem}{B\arabic{problem}}
\theoremstyle{remark}
\newtheorem*{remark}{Remark}
\newenvironment{answer}
  {\par\medskip\noindent\textbf{Answer.}\begin{quotation}\noindent}
  {\end{quotation}\medskip}

% ---------------------------------------------------------------------------
% Grader margin notes. Rendered form of the `% GRADER[...][...][...]:` comment
% convention (see template_answers.tex). Use inline where a visible note helps:
%   \gradernote{B2/induction-soundness/2: the invariant is not preserved by the last assignment}
% ---------------------------------------------------------------------------
\definecolor{gradercolor}{RGB}{160, 30, 30}
\newcommand{\gradernote}[1]{\textcolor{gradercolor}{\small\sffamily [GRADER: #1]}}

% ---------------------------------------------------------------------------
% Code listings: Dafny, Verus (Rust inside verus!{...}), Python.
% ---------------------------------------------------------------------------
\definecolor{codekeyword}{RGB}{0, 60, 130}
\definecolor{codecomment}{RGB}{90, 110, 90}
\definecolor{codestring}{RGB}{130, 60, 10}
\definecolor{codehole}{RGB}{180, 0, 0}

\lstdefinelanguage{dafny}{
  sensitive=true,
  morekeywords={
    abstract, allocated, as, assert, assume, bool, break, by, calc, case, char,
    class, codatatype, const, constructor, continue, datatype, decreases, else,
    ensures, exists, expect, export, extends, false, forall, fresh, function,
    ghost, if, import, in, include, int, invariant, is, iterator, label, lemma,
    map, match, method, modifies, modify, module, multiset, nat, new, newtype,
    null, old, opaque, opened, predicate, print, provides, reads, real, refines,
    requires, return, returns, reveal, reveals, seq, set, static, string, then,
    this, trait, true, twostate, type, unchanged, var, while, witness},
  emph={assume},
  emphstyle=\color{codehole}\bfseries,
  comment=[l]{//},
  morecomment=[s]{/*}{*/},
  morestring=[b]",
}

\lstdefinelanguage{verus}{
  sensitive=true,
  morekeywords={
    as, break, const, continue, crate, dyn, else, enum, false, fn, for, if,
    impl, in, let, loop, match, mod, move, mut, pub, ref, return, self, Self,
    static, struct, super, trait, true, type, unsafe, use, where, while,
    spec, proof, exec, open, closed, tracked, ghost, requires, ensures,
    recommends, invariant, decreases, assert, forall, exists, choose, by,
    broadcast, uninterp, axiom, returns, int, nat, Seq, Set, Map, Ghost, Tracked},
  emph={assume, admit},
  emphstyle=\color{codehole}\bfseries,
  comment=[l]{//},
  morecomment=[s]{/*}{*/},
  morestring=[b]",
}

\lstset{
  basicstyle=\ttfamily\small,
  keywordstyle=\color{codekeyword}\bfseries,
  commentstyle=\color{codecomment}\itshape,
  stringstyle=\color{codestring},
  showstringspaces=false,
  columns=fullflexible,
  keepspaces=true,
  breaklines=true,
  frame=leftline,
  framerule=0.6pt,
  xleftmargin=1em,
  aboveskip=0.8em,
  belowskip=0.8em,
  literate=
    {→}{{$\to$}}1        {←}{{$\leftarrow$}}1  {⇒}{{$\Rightarrow$}}1
    {⇔}{{$\Leftrightarrow$}}1                  {↦}{{$\mapsto$}}1
    {∀}{{$\forall$}}1    {∃}{{$\exists$}}1     {λ}{{$\lambda$}}1
    {∧}{{$\wedge$}}1     {∨}{{$\vee$}}1        {¬}{{$\neg$}}1
    {≠}{{$\neq$}}1       {≤}{{$\leq$}}1        {≥}{{$\geq$}}1
    {∈}{{$\in$}}1        {∉}{{$\notin$}}1      {⊆}{{$\subseteq$}}1
    {×}{{$\times$}}1     {·}{{$\cdot$}}1       {∅}{{$\varnothing$}}1
    {ℕ}{{$\mathbb{N}$}}1 {ℤ}{{$\mathbb{Z}$}}1
    {⊢}{{$\vdash$}}1     {⟦}{{$\llbracket$}}1  {⟧}{{$\rrbracket$}}1,
}

% Code environments and inline code.
\lstnewenvironment{dafnycode}{\lstset{language=dafny}}{}
\lstnewenvironment{veruscode}{\lstset{language=verus}}{}
\lstnewenvironment{pycode}{\lstset{language=Python}}{}
\newcommand{\dfy}[1]{\lstinline[language=dafny]{#1}}
\newcommand{\vrs}[1]{\lstinline[language=verus]{#1}}
\newcommand{\py}[1]{\lstinline[language=Python]{#1}}


\title{Q03 --- First syntax II: modes, logic, tooling, trust \\ \large Written problems}
\author{dafny-verus-learn-0}
\date{}

\begin{document}
\maketitle

\section*{Instructions}
Answer in \texttt{Q03\_answers.tex}, one section per problem; keep its
\verb+\label{ans:B<n>}+ anchors. \texttt{make -C paper paper} must stay green:
a non-compiling answer set is the written track's red build.

Order: B1 and B2 after chapter section~2 (the S8 lab, 20 minutes); B3 and B4
before you start the holes (15 minutes). Write B3 \emph{before} you open the
hints under H1 in the exercise file. Then run \texttt{/grade-paper Q03} in a
fresh session. Once H1d verifies, \verb+/grade-paper Q03 --post+ re-scores
B3's blueprint against your own proof.

Hand calculations are graded step by step: write every step, and name the rule
it uses. Derive first; check with the verifier afterwards, as chapter section~2
does. Compare with \texttt{paper/solutions/} only as the module's last step.

Notation: $\WP{c}{Q}$ is the weakest precondition of program $c$ for
postcondition $Q$; $\subst{Q}{x}{e}$ is $Q$ with $e$ substituted for $x$.
All variables range over the mathematical integers unless a type says
otherwise.

% ---------------------------------------------------------------------------
\begin{problem}[hand-calculation]\label{prob:B1}
\textbf{The weakest precondition of loop-free code} (10 minutes). Over integer
variables $x$ and $y$, the program $c$ is
\[
  y \assign x + 3 \seqc
  \mathbf{if}\ y > 10\ \mathbf{then}\ y \assign y - 10\ \mathbf{else}\ y \assign 2 \cdot y \seqc
  x \assign y - 1
\]
and the postcondition is $Q \equiv x \geq 4$. In the two languages, with the
precondition $P \equiv 0 \leq x \leq 8$ of part (b) attached:
\begin{dafnycode}
method B1(x0: int) returns (x: int, y: int)
  requires 0 <= x0 <= 8
  ensures x >= 4
{
  x := x0;
  y := x + 3;
  if y > 10 { y := y - 10; } else { y := 2 * y; }
  x := y - 1;
}
\end{dafnycode}
\begin{veruscode}
proof fn b1(x0: int) -> (x: int)
    requires 0 <= x0 <= 8,
    ensures x >= 4,
{
    let y = x0 + 3;
    let y = if y > 10 { y - 10 } else { 2 * y };
    y - 1
}
\end{veruscode}
\begin{enumerate}
\item[(a)] Compute $\WP{c}{Q}$ backwards, one statement at a time, naming the
  rule at each step (assignment, sequence, conditional). Then simplify the
  result to a formula over the start value of $x$ alone, justifying each
  simplification over the integers.
\item[(b)] State the verification condition $P \impl \WP{c}{Q}$ and decide
  it. If it fails, exhibit a start state that satisfies $P$ and ends outside
  $Q$, and trace the program on it.
\item[(c)] In one sentence: is the set of start values of $x$ that satisfy
  $\WP{c}{Q}$ an interval?
\end{enumerate}
\end{problem}

% ---------------------------------------------------------------------------
\begin{problem}[hand-calculation]\label{prob:B2}
\textbf{Verification conditions of a loop} (10 minutes). The loop below
counts $x$ up to $n$ while it counts $y$ down from $n$:
\begin{dafnycode}
method Countdown(n: int) returns (x: int, y: int)
  requires n >= 0
  ensures y == 0
{
  x, y := 0, n;
  while x < n
    invariant x <= n && x + y == n
    decreases n - x
  {
    x := x + 1;
    y := y - 1;
  }
}
\end{dafnycode}
\begin{veruscode}
fn countdown(n: u64) -> (y: u64)
    ensures y == 0,
{
    let mut x: u64 = 0;
    let mut y: u64 = n;
    while x < n
        invariant
            x <= n,
            x + y == n,
        decreases n - x,
    {
        x = x + 1;
        y = y - 1;
    }
    y
}
\end{veruscode}
Work with the Dafny version (integers) in (a) and (b). Write $I$ for the
invariant $x \leq n \wedge x + y = n$, $b$ for the guard $x < n$ and $S$ for
the body.
\begin{enumerate}
\item[(a)] Write the four verification conditions of the loop as formulas:
  \emph{entry} (the precondition and the initialisation establish $I$),
  \emph{maintenance} ($I \wedge b \impl \WP{S}{I}$, with $\WP{S}{I}$ computed by
  the assignment and sequence rules), \emph{use at exit} ($I \wedge \neg b$
  gives the postcondition) and \emph{termination} (the measure
  $n - x$ is bounded below while the loop runs, and $S$ decreases it).
  Decide each with a one-line argument.
\item[(b)] Weaken the invariant to $I' \equiv x + y = n$ and redo the four
  conditions. Which one fails? Exhibit a state that satisfies its hypotheses
  and violates its conclusion. Can any run of the program reach that state?
  What does the failure say about the program, and what about $I'$?
\item[(c)] Keep $I'$, but give $y$ a type that cannot go below zero: \dfy{y: nat}
  in Dafny, or the \vrs{u64} of the Verus version. Does the loop verify now?
  Name the fact that changes the outcome of the condition that failed in (b),
  and say where the verifier gets it.
\end{enumerate}
\end{problem}

% ---------------------------------------------------------------------------
\begin{problem}[informal-proof-first]\label{prob:B3}
\textbf{A blueprint for H1d} (7 minutes). Write the informal proof of the
lemma below \emph{before} you formalize it, and before you read the hints
under it in the exercise file. The definitions are those of the given layers
L1 and L2:
\begin{dafnycode}
const BASE: nat := 18446744073709551616

predicate Valid(s: seq<nat>)
{
  forall i | 0 <= i < |s| :: s[i] < BASE
}

function ToNat(s: seq<nat>): nat
  decreases |s|
{
  if |s| == 0 then 0 else s[0] + BASE * ToNat(s[1..])
}

function AddLimb(s: seq<nat>, d: nat): seq<nat>
  decreases |s|
{
  if |s| == 0 then (if d == 0 then [] else [d])
  else if s[0] + d < BASE then [s[0] + d] + s[1..]
  else [s[0] + d - BASE] + AddLimb(s[1..], 1)
}

lemma AddLimbToNat(s: seq<nat>, d: nat)
  requires Valid(s) && d < BASE
  ensures ToNat(AddLimb(s, d)) == ToNat(s) + d
  decreases |s|
\end{dafnycode}
State the induction you use and the induction hypothesis as a formula, cover
every case, and name every helper fact by its statement: a lemma of L2 (read
its file) or a new lemma you would have to prove.

\emph{What will be formalized:} Q03.H1d, \dfy{AddLimbToNat} in the Dafny
exercise file
\begin{quote}
\texttt{dafny/Q/Q03\_FirstSyntaxII/core/Q03\_FirstSyntaxII.dfy}.
\end{quote}
\noindent That unit does
no induction for you: each use of the induction hypothesis is a recursive call
you write, and each helper fact is a lemma call or an \dfy{assert}. Your proof
is the blueprint: someone else should be able to formalize it without
inventing any statement. (The Verus pair H1v rests on the given
\vrs{lemma_add_limb_lift}, whose per-limb relation is the arithmetic of the
carry case.)
\end{problem}

% ---------------------------------------------------------------------------
\begin{problem}[paper-card]\label{prob:B4}
\textbf{Paper card 02.13, Clover: Closed-Loop Verifiable Code Generation}
(8 minutes). Fill the card on its eight fields: task shape, oracle, what is
trusted, cheat surface, metric, budget, variance, artifact. Each field is
graded against the pinned range of your own \texttt{cheatsheets.md}: lines
944--967, the entry from its heading to its \emph{Why it matters} paragraph,
in the copy the command below names. The pin is the sha256 of exactly those
lines, computed with the first line below; it must print the second.
{\small
\begin{verbatim}
sed -n '944,967p' ~/code/research/agentic-formal-veri/cheatsheets.md | sha256sum
d323712a0a7d430510ca9b9f14daa83c0c6ca9ad876eaffa542ed9c6bc56a38e  -
\end{verbatim}
}
Write each field in your own words, a sentence or two. Where the range says
nothing about a field, write that, and what you would need to know. A hash
that no longer matches means the card file changed: the grader refuses the
card until the pin is renewed.
\end{problem}

\end{document}
